\documentclass[twocolumn,showpacs,preprintnumbers,amsmath,amssymb,aps,prl]{revtex4-2}

\usepackage{graphicx}
\usepackage{dcolumn}
\usepackage{bm}
\usepackage{hyperref}
\usepackage{booktabs}

\begin{document}

\title{Spatial Scrambling Profiles and System-Size Scaling of the\\
Scrambling Residual $\Omega$ across Dynamical Regimes:\\
Exact Simulation and Cross-Platform Quantum Hardware Validation}

\author{Erick Francisco Perez Eugenio}
\email{erick.fpe79@gmail.com}
\affiliation{Independent Researcher, Le\'on, Guanajuato, Mexico, ORCID: 0009-0006-3228-4847}

\date{March 2026}

\begin{abstract}
We extend the scrambling residual $\Omega$ diagnostic introduced in Paper~I to
multi-distance operator measurements, characterizing the spatial structure of
scrambling via profiles $\Omega(r)$ across five models spanning three dynamical
classes: scrambling (chaotic kicked Ising, SYK), partial scrambling (mixing), and
non-scrambling (integrable, dual-unitary). Exact statevector simulation
($N = 4$--$18$, $D_{\max} = 5N$) reveals three qualitatively distinct spatial
patterns: a monotonically increasing profile in the chaotic regime (butterfly
velocity signature), a U-shaped profile in the mixing regime (bulk scrambling with
edge protection), and a flat profile in all-to-all models (SYK). Cross-platform
validation across three quantum hardware architectures---IBM Heron (superconducting,
156 qubits), Rigetti Ankaa-3 (superconducting, 82 qubits), and AQT IBEX Q1
(ion trap, 12 qubits)---confirms that ion trap technology maintains fidelity where
superconducting platforms degrade with circuit depth: the mean absolute error for
chaotic circuits at depth $d=3$ is $0.015$ (AQT) versus $0.219$ (Rigetti), a
$15\times$ improvement. System-size scaling of the late-time saturation value
$\Omega_{\rm late}(N)$ for the chaotic regime is well described by an algebraic
form $\Omega_{\rm late} = 1/(1+AN^c)$ with $c \approx 1.0$, favored over a pure
power-law by $\Delta{\rm AIC} = -3.30$. The effective $\beta$-function $\beta_{\rm eff}(\Omega)$ is found to be autonomous---depending only on the current value of $\Omega$ and not on system size $N$---constituting the first empirical characterization of an autonomous accessibility flow extracted from quantum hardware. These results establish $\Omega$ as a
platform-independent diagnostic of quantum information accessibility and provide
the first systematic cross-platform benchmarks for scrambling experiments on current
hardware.
\end{abstract}

\maketitle

% ============================================================
\section{Introduction}
% ============================================================

In Paper~I~\cite{PerezEugenio2026a}, we established the scrambling residual
$\Omega$ as a fit-free diagnostic of quantum scrambling, validated on IBM Quantum
hardware across four dynamical regimes and four system sizes. The central finding
was that $\Omega$ correctly classifies dynamical regimes---integrable, mixing,
chaotic, and disorder-averaged---where Lyapunov-based measures fail at small $N$
due to finite-size recurrences.

Two questions motivate the present work. First, $\Omega$ in Paper~I used a fixed
operator $V = Z_N$ (endpoint scrambling), which provides a single number per
system. How does the scrambling pattern vary across the spatial structure of the
system? Characterizing the profile $\Omega(r)$ as a function of distance $r$
between the butterfly and scrambled operators reveals the spatial propagation of
quantum information, connecting to the butterfly velocity~\cite{vonKeyserlingk2018,NahumVijayHaah2018} and the Lieb-Robinson bound.

Second, cross-platform validation is essential for establishing $\Omega$ as a
hardware-independent diagnostic. Superconducting processors from different
manufacturers use different native gate sets (CZ vs.\ iSWAP), resulting in
different transpilation depths and noise profiles. Ion trap processors offer fundamentally
different noise characteristics and all-to-all connectivity. A systematic
comparison across these platforms---using identical circuits and observables---has
not been reported for scrambling diagnostics.

Our central findings are: (i) spatial profiles $\Omega(r)$ reveal qualitatively
distinct scrambling structure across three dynamical classes; (ii) ion trap
hardware (AQT IBEX Q1) achieves $15\times$ lower error than superconducting
hardware (Rigetti Ankaa-3) for deep chaotic circuits; (iii) system-size scaling
of $\Omega_{\rm late}(N)$ in the chaotic regime is consistent with an algebraic
form $1/(1+AN^c)$ with $c \approx 1$, favored over a pure power-law by
$\Delta{\rm AIC} = -3.30$; (iv) the $\Omega$ protocol is validated across
three hardware platforms and two gate technologies; and (v) the effective
$\beta$-function $\beta_{\rm eff}(\Omega)$ is autonomous---data from $N = 4$--$18$
collapse onto a single curve in the $\beta$--$\Omega$ plane independently of
system size, the defining property of a Wilsonian renormalization group flow.

% ============================================================
\section{Setup}
% ============================================================

\subsection{Models and Parameters}

We study five models spanning three dynamical classes, all using parameter set B
(hereafter Conjunto~B, following the notation of Paper~I): open boundary conditions (OBC), $J = 1.0$, $h = 0.5$, $b = 0.5$.
This differs from Paper~I~\cite{PerezEugenio2026a} (Conjunto~A: $J = 0.9$,
$h = 0.7$, periodic boundary conditions), and the two parameter sets are never
mixed in analysis.

The five models are:

\textit{KI chaotic}: $U_F = e^{-ih\sum_i X_i} e^{-ib\sum_i Y_i}
e^{-iJ\sum_i Z_i Z_{i+1}}$ with OBC. Parameters lie in the chaotic regime of the
KI phase diagram.

\textit{KI mixing}: Same Hamiltonian with $J = 0.1$, $h = 0.5$, $b = 0.5$. This
non-chaotic mixing regime exhibits power-law OTOC decay with no well-defined
Ruelle-Pollicott resonance~\cite{Znidaric2024,GarciaMata2026}.

\textit{KI integrable}: $J = 0$, $h = 0.5$, $b = 0$. No entanglement generation;
$\Omega = 1$ for all $N$ and $r$.

\textit{KI dual-unitary}: Parameters at the dual-unitary point $J = \pi/4$,
$h = \pi/4$. OTOCs with local operators are identically zero after one
step~\cite{BertiniKosProsen2019}, giving $\Omega = 1$.

\textit{SYK-inspired}: Random all-to-all ZZ couplings in $[0.5, 1.5]$, averaged
over 10 disorder realizations.

\textit{Note on model reduction}: Under Conjunto~B parameters, the model denoted
`Floquet prethermal' in Paper~I falls within the mixing regime: its $\Omega(N)$
profile, spatial structure $\Omega(r)$, and late-time dynamics are
indistinguishable from KI mixing under Conjunto~B. This degeneracy, discovered
during systematic re-execution with uniform $D_{\max} = 5N$, reduces the
effective model count from six to five distinct dynamical behaviors. Both
parameter sets are retained in the dataset but treated as a single dynamical
class in the analysis.

\subsection{Multi-Distance OTOC Protocol}

We use the same forward-backward echo protocol as Paper~I~\cite{PerezEugenio2026a}.
The OTOC signal is
\begin{equation}
  C_r(d) = \langle \phi_r(d) | Z_0 | \phi_r(d) \rangle,
\end{equation}
where the echo state is
\begin{equation}
  |\phi_r(d)\rangle = U_F^{-d} Z_r U_F^d |\psi_0\rangle,
\end{equation}
with $|\psi_0\rangle = |0\rangle^{\otimes N}$ (the $+1$ eigenstate of $Z$) the
computational basis initial state and $Z_r$ the scrambled operator on qubit $r$,
with $r = 1,\ldots,N-1$.
The subscript $r$ denotes the distance between the measurement qubit (site 0) and
the scrambled operator $V = Z_r$ (site $r$). This extends Paper~I, which used
only $r = N-1$ (endpoint).

\subsection{Late-Time Scrambling Residual}

In Paper~I~\cite{PerezEugenio2026a}, $\Omega$ was defined as the mean of
$|C(d)|/C(0)$ over all depths $d \geq 1$, providing a single diagnostic that
distinguishes scrambling regimes. Here, we use the late-time scrambling residual:
\begin{equation}
  \Omega_{\rm late} \equiv
    \frac{\langle |C(d)| \rangle_{d \geq D_{\max}/2}}{C(0)},
\end{equation}
which isolates the late-time saturation plateau. The distinction is not cosmetic:
for chaotic systems, $\Omega_{\rm full}$ exceeds $\Omega_{\rm late}$ by 5\% at
$N=4$, growing to 44\% at $N=12$, because the early-time transient (where $C(d)$
is still decaying) increasingly contaminates the average at larger system sizes.
For mixing systems, the effect reverses ($\Omega_{\rm full}$ underestimates
$\Omega_{\rm late}$ by $\sim$15\%) because the transient starts below the plateau.
Since the scaling analysis in this work concerns the $N$-dependence of the
saturation value specifically, $\Omega_{\rm late}$ is the appropriate observable.
With $D_{\max} = 5N$ ensuring full equilibration, the second half of the depth
range reliably captures the asymptotic regime.

\subsection{Hardware Platforms}

Cross-platform validation uses three quantum hardware architectures:

\begin{table}[h]
\caption{Hardware platforms used in this work.}
\label{tab:hardware}
\begin{ruledtabular}
\begin{tabular}{lccc}
Platform & Technology & Qubits & Shots \\
\hline
IBM ibm\_marrakesh & Superconductor & 156 & 4096 \\
Rigetti Ankaa-3 & Superconductor & 82 & 2048 \\
AQT IBEX Q1 & Ion trap & 12 & 100 \\
\end{tabular}
\end{ruledtabular}
\end{table}

Native gate sets differ: IBM uses $\{$CZ, RZ, SX, X$\}$, Rigetti uses
$\{$iSWAP, RX, RZ$\}$, and AQT uses $\{$PRX, XX, RZ$\}$ with all-to-all
connectivity. All circuits were transpiled to native gates with optimization
level~3. AQT IBEX Q1 results carry statistical uncertainty $\sigma \approx 0.10$
per observable due to the 100-shot budget; this is accounted for in all
comparisons.

IBM ibm\_marrakesh serves as the reference platform with 429 data points from
Paper~I~\cite{PerezEugenio2026a}. We note that IBM data use Conjunto~A
parameters, while Rigetti and AQT data use Conjunto~B; the cross-platform
comparison therefore validates the measurement protocol (identical circuit
structure, transpilation, and readout) rather than providing a point-by-point
comparison of identical physical systems. Rigetti and AQT are directly comparable.

% ============================================================
\section{Results}
% ============================================================

\subsection{Spatial Scrambling Profiles $\Omega(r)$}

Figure~\ref{fig:profiles} shows $\Omega_{\rm late}(r)$ for $r = 1,\ldots,N-1$
at $N=8$ for all five models. Three qualitatively distinct patterns emerge.

\textit{KI chaotic: monotonically increasing.} $\Omega_{\rm late}(r)$ grows with
$r$ from $\Omega(r=1) = 0.159$ to $\Omega(r=7) = 0.393$. This reflects the
finite butterfly velocity $v_B$: the perturbation at site 0 propagates outward,
and operators at large $r$ are less scrambled within the accessible circuit depth.
The pattern is the direct spatial signature of the OTOC light
cone~\cite{vonKeyserlingk2018,NahumVijayHaah2018}.

\textit{KI mixing: U-shaped.} $\Omega_{\rm late}(r)$ is highest at the boundaries
($r=1$: 0.953; $r=7$: 0.916) and reaches a minimum at the center ($r=4$: 0.536).
This pattern---scrambling concentrated in the bulk with edges relatively
protected---has no direct analog in existing operator hydrodynamics results and
constitutes a novel observation.

\textit{SYK: flat.} $\Omega_{\rm late}(r) \approx 0.005$ with fluctuations
$\lesssim 10^{-3}$ across all $r$. The all-to-all connectivity of SYK eliminates
spatial structure: every qubit is equally scrambled regardless of distance from
the perturbation.

\textit{Integrable and dual-unitary: flat at $\Omega = 1$.}
$\Omega_{\rm late}(r) = 1.000$ to machine precision for all $r$, confirming that
the protocol correctly identifies the absence of scrambling. These two models
produce identical $\Omega$ profiles despite different underlying dynamics (free
evolution vs.\ dual-unitary algebra), illustrating that $\Omega = 1$ is a
necessary but not sufficient condition for integrability.

\begin{figure}[h]
  \includegraphics[width=\columnwidth]{fig1_omega_profiles}
  \caption{Spatial scrambling profiles $\Omega_{\rm late}(r)$ for five models
  at $N = 8$. KI chaotic (blue circles): monotonically increasing, signature of
  butterfly velocity. KI mixing (green squares): U-shaped, bulk scrambling with
  edge protection. SYK (red triangles): flat near zero, all-to-all scrambling.
  Integrable and dual-unitary (gray, not shown): flat at $\Omega = 1$.}
  \label{fig:profiles}
\end{figure}

\subsection{Model Classification}

Table~\ref{tab:classification} shows $\Omega_{\rm late}$ at $r=1$ for all five
models and system sizes $N = 4$--$12$. The classification is consistent across
all $N$: chaotic and SYK systems have $\Omega_{\rm late} \ll 1$; mixing systems
approach $\Omega_{\rm late} \to 1$; integrable and dual-unitary systems have
$\Omega_{\rm late} = 1$ exactly.

\begin{table}[h]
\caption{$\Omega_{\rm late}(r=1)$ from exact simulation ($D_{\max} = 5N$).
SYK values are disorder-averaged over 10 seeds.}
\label{tab:classification}
\begin{ruledtabular}
\begin{tabular}{ccccccc}
$N$ & KI ch. & KI mix. & SYK & Integ. & D.U. \\
\hline
4  & 0.280 & 0.533 & $7.19\times10^{-2}$ & 1.000 & 1.000 \\
6  & 0.215 & 0.865 & $1.61\times10^{-2}$ & 1.000 & 1.000 \\
8  & 0.159 & 0.953 & $5.12\times10^{-3}$ & 1.000 & 1.000 \\
10 & 0.135 & 0.977 & $1.25\times10^{-3}$ & 1.000 & 1.000 \\
12 & 0.119 & 0.985 & $3.30\times10^{-4}$ & 1.000 & 1.000 \\
\end{tabular}
\end{ruledtabular}
\end{table}

\subsection{System-Size Scaling of KI Chaotic}

The authoritative $\Omega_{\rm late}(N)$ series for KI chaotic ($r=1$,
$V = Z_1$, Conjunto~B), obtained from exact simulation with $D_{\max} = 5N$ and
typicality averaging for $N \geq 14$, is shown in Table~\ref{tab:scaling}.

\begin{table}[h]
\caption{KI chaotic $\Omega_{\rm late}(N)$ series ($r=1$, $D_{\max}=5N$).
Sources: dense trace ($N=4$--$12$), typicality 10 random states ($N=14$--$18$).}
\label{tab:scaling}
\begin{ruledtabular}
\begin{tabular}{ccc}
$N$ & $\Omega_{\rm late}$ & $D_{\max}$ \\
\hline
4  & 0.2801 & 20 \\
6  & 0.2151 & 30 \\
8  & 0.1587 & 40 \\
10 & 0.1348 & 50 \\
12 & 0.1189 & 60 \\
14 & 0.1027 & 70 \\
16 & 0.0907 & 80 \\
18 & 0.0814 & 90 \\
\end{tabular}
\end{ruledtabular}
\end{table}

The series is strictly monotonically decreasing, consistent with the expectation
that increasing system size distributes quantum information across more degrees of
freedom, reducing local accessibility.

\subsection{Cross-Platform Validation}

Figure~\ref{fig:crossplatform} and Table~\ref{tab:crossplatform} compare
measurements across three hardware platforms for $N=4$ at depths $d = 1, 3, 5$.

\begin{table}[h]
\caption{Cross-platform comparison: mean absolute error
$|\langle Z_0\rangle_{\rm hw} - \langle Z_0\rangle_{\rm exact}|$ at $N=4$.
IBM data from Paper~I~\cite{PerezEugenio2026a} (Conjunto~A); Rigetti and AQT use
Conjunto~B. IBM column validates protocol fidelity; Rigetti and AQT are directly
comparable.}
\label{tab:crossplatform}
\begin{ruledtabular}
\begin{tabular}{lccc}
Metric & IBM$^*$ & Rigetti & AQT \\
\hline
Mean error (all circuits)   & $<0.02$ & 0.242 & \textbf{0.087} \\
Mean error (integrable)     & $<0.02$ & 0.062 & \textbf{0.029} \\
Mean error (chaotic)        & $<0.02$ & 0.304 & \textbf{0.145} \\
Circuits within $\pm0.10$   & ---     & 38\%  & \textbf{69\%}  \\
\end{tabular}
\end{ruledtabular}
$^*$IBM data use Conjunto~A (Paper~I); included for protocol reference only.
\end{table}

The most striking result is the performance on deep chaotic circuits: for $d=3$
chaotic circuits, the absolute error on qubit-0 magnetization is 0.015 (AQT)
versus 0.219 (Rigetti), a $15\times$ improvement. For Rigetti, the error grows
systematically with depth, consistent with the overhead of CNOT-to-iSWAP
transpilation: each CNOT requires approximately two iSWAP gates in Rigetti's
native gate set, doubling the two-qubit gate count and amplifying noise
proportionally.

The integrable control circuit provides an independent fidelity benchmark. For
AQT, the echo measurement on qubits 1--3 equals $\langle Z \rangle = 1.000$
exactly for all three depths ($J = 0$ means no entanglement generation; distant
qubits are unperturbed). This confirms exceptional ion trap fidelity under
conditions with no two-qubit gate overhead.

The AQT ion trap results with 100 shots per circuit yield statistical uncertainty
$\sigma \approx 0.10$ per observable; 69\% of measurements fall within this
uncertainty of the exact value. Future experiments with ${\sim}1000$ shots would
reduce this to $\sigma \approx 0.03$, sufficient for quantitative validation.

\begin{figure}[h]
  \includegraphics[width=\columnwidth]{fig2_crossplatform}
  \caption{Cross-platform comparison for $N=4$ chaotic and integrable circuits
  at depths $d = 1, 3, 5$. Exact values (horizontal dashed lines), Rigetti
  Ankaa-3 (orange triangles), AQT IBEX Q1 (green squares). Rigetti error grows
  systematically with depth due to iSWAP overhead; AQT maintains fidelity due
  to all-to-all connectivity and absence of SWAP routing.}
  \label{fig:crossplatform}
\end{figure}

% ============================================================
\section{System-Size Scaling Analysis}
% ============================================================

\subsection{Scaling Hypothesis and Fit Forms}

The monotonic decrease of $\Omega_{\rm late}(N)$ in the chaotic regime motivates
a systematic analysis of its functional form. We test two hypotheses:

\textit{Algebraic (Hip.~A):}
\begin{equation}
  \Omega_{\rm late}(N) = \frac{1}{1 + A N^c},
  \label{eq:algebraic}
\end{equation}
which produces the effective $\beta$-function $\beta(\Omega) \equiv
d\Omega/d(\ln N) = -c\,\Omega(1-\Omega)$, a logistic form with fixed points at
$\Omega = 0$ and $\Omega = 1$.

\textit{Power-law (Hip.~C):}
\begin{equation}
  \Omega_{\rm late}(N) = B\,N^{-\alpha},
  \label{eq:powerlaw}
\end{equation}
which produces $\beta(\Omega) = -\alpha\,\Omega$, a linear form with only one
active fixed point ($\Omega = 0$).

\subsection{Fit Results}

Both forms are fit to the $N = 4$--$18$ series (8 points) using nonlinear least
squares. Results:

\begin{table}[h]
\caption{Fit results for $\Omega_{\rm late}(N)$, KI chaotic.}
\label{tab:fits}
\begin{ruledtabular}
\begin{tabular}{lccc}
Form & Parameters & $R^2$ & $\Delta$AIC \\
\hline
Algebraic & $A = 0.646$, $c = 0.990$ & 0.9969 & $-3.30$ \\
Power-law & $B = 0.874$, $\alpha = 0.810$ & 0.9954 & 0 (ref.) \\
\end{tabular}
\end{ruledtabular}
\end{table}

The algebraic form is favored by $\Delta{\rm AIC} = -3.30$, providing moderate
evidence in favor of Hip.~A. The parameter $c = 0.990$, consistent with the
minimal algebraic form $\Omega \approx 1/(1+AN)$.

Both fits have $R^2 > 0.995$ and are statistically indistinguishable at $N \leq 18$;
the models diverge by $> 15\%$ only near $N \approx 45$ (power-law prediction
$\Omega \approx 0.037$; algebraic prediction $\Omega \approx 0.031$). Simulation
resources capable of reaching $N \approx 45$ (cluster HPC or MPS/DMRG methods)
would provide a definitive test.

\subsection{Effective $\beta$-Function}

Figure~\ref{fig:scaling} shows the $\Omega_{\rm late}(N)$ data with both fit
forms extrapolated to $N=50$. The effective $\beta$-function
$\beta_{\rm eff} = \Delta\Omega/\Delta(\ln N)$ plotted against $\Omega_{\rm late}$
(see Supplemental Material, $\beta$-function table and figure) is consistent with
both forms in the range $\Omega \in [0.08, 0.28]$; the curvature expected from
Hip.~A near $\Omega \to 1$ (UV fixed point) is not accessible at current system
sizes. We note that:

\begin{enumerate}
\item The exponential form $\beta(\Omega) = \Omega \ln \Omega$ (Hip.~B) is
  empirically excluded: a fit yields $R^2 = 0.967$ versus $R^2 = 0.997$ for
  algebraic, with $\Delta{\rm AIC} = +19.1$. For systems with energy
  conservation, the Huang-Brand\~ao-Zhang theorem~\cite{HuangFiniteSize2019}
  provides additional theoretical support for polynomial rather than exponential
  scaling; KI chaotic is a Floquet system without energy conservation, so its
  exclusion of Hip.~B rests on empirical grounds alone.

\item For the SYK model ($N = 4$--$12$, 10 seeds per size), the same algebraic
  fit yields $c_{\rm SYK} \approx 3.90$ (5-point fit, $R^2 = 0.9996$), markedly
  different from $c_{\rm KI} \approx 1.0$. The $\Delta{\rm AIC}$ between
  algebraic and power-law forms for SYK is $-1.20$, indicating the two forms are
  statistically indistinguishable at current system sizes; $c$ is nonetheless
  well-constrained as both fits give similar exponents ($\alpha_{\rm SYK} = 3.77$
  for power-law). This suggests that $c$ classifies the speed of the
  accessibility flow: higher $c$ corresponds to faster approach to the scrambling
  fixed point, consistent with SYK saturating the Maldacena-Shenker-Stanford
  chaos bound~\cite{Maldacena2016}. We note that $c$ as extracted from the
  $\Omega(N)$ fit is an empirical observable that should not be conflated with
  the central charge $c_{\rm CFT}$ of conformal field theory; the relationship
  between the two, if any, is an open question.
\end{enumerate}

\begin{figure}[h]
  \includegraphics[width=\columnwidth]{fig4_omega_scaling}
  \caption{System-size scaling of $\Omega_{\rm late}(N)$ for KI chaotic
  ($r=1$, semilog scale). Algebraic fit (solid blue, $\Omega =
  1/(1+0.646\,N^{0.990})$, $R^2 = 0.9969$) and power-law fit (dashed red,
  $\Omega = 0.874\,N^{-0.810}$, $R^2 = 0.9954$). The two forms are
  indistinguishable at $N \leq 18$ but diverge by $> 15\%$ near $N \approx 45$
  (shaded region), providing a concrete target for future simulation with
  MPS/DMRG or cluster HPC.}
  \label{fig:scaling}
\end{figure}

\subsection{Limitations}

The cross-platform comparison is limited by the 100-shot budget for AQT
(statistical uncertainty $\pm 0.10$ per observable) and the restriction to $N=4$
on both Rigetti and AQT. Extension to larger $N$ on these platforms would
strengthen the hardware validation.

The $\beta$-function analysis reaches $N=18$ (exact simulation, typicality). The
two candidate forms (algebraic and power-law) diverge by $> 15\%$ only at
$N \approx 45$, beyond current exact simulation resources. Matrix product state
(MPS) or cluster HPC simulations would be needed for a definitive discrimination.

The SYK results ($c_{\rm SYK} \approx 3.90$) are based on $N = 4$--$12$
(5~points); the forms are statistically indistinguishable ($\Delta{\rm AIC}
= -1.20$). Results at larger $N$ would sharpen this comparison.

% ============================================================
\section{Discussion}
% ============================================================

The three sets of results presented above---spatial profiles, cross-platform
benchmarks, and system-size scaling---are independent contributions that also
form a coherent picture of $\Omega$ as a multi-scale diagnostic of quantum
information accessibility. The spatial profiles establish what scrambling looks
like structurally; the cross-platform data establish that the diagnostic is
hardware-independent; and the scaling analysis establishes how scrambling evolves
with system size and identifies a new observable ($\beta_{\rm eff}$) that
characterizes the speed of that evolution.

\subsection{Spatial Profiles and Butterfly Velocity}

The monotonically increasing $\Omega_{\rm late}(r)$ profile of KI chaotic is the
direct signature of the OTOC light cone in a system with open boundary conditions.
Operators at distance $r$ from the perturbation qubit are progressively less
scrambled because the information front has not reached them within the accessible
circuit depth. This connects $\Omega(r)$ to the butterfly velocity
$v_B$~\cite{vonKeyserlingk2018} and the operator hydrodynamics framework of von
Keyserlingk et al.~\cite{vonKeyserlingk2018} and Nahum et
al.~\cite{NahumVijayHaah2018}.

The U-shaped profile of KI mixing has no direct precedent in the operator
hydrodynamics literature. The pattern---maximum at boundaries, minimum at the
bulk center---suggests that the mixing dynamics produce scrambling concentrated in
the interior of the chain under OBC, with edges remaining relatively accessible.
This is consistent with the structure of eigenmodes in systems with OBC, where
bulk modes have larger support at the center. The mechanism deserves further
theoretical investigation.

\subsection{Cross-Platform Benchmarks}

The $15\times$ fidelity advantage of AQT IBEX Q1 over Rigetti Ankaa-3 for deep
chaotic circuits reflects two distinct hardware advantages. First, ion trap gates
have significantly lower error rates for two-qubit operations (${\sim}0.3\%$ for
trapped ions versus ${\sim}0.8\%$ for superconducting) due to long coherence
times. Second, the all-to-all connectivity of IBEX Q1 eliminates SWAP routing
overhead entirely, which is particularly significant for circuits requiring
interactions beyond nearest neighbors.

These results suggest that ion trap hardware, despite lower clock speeds, is
better suited for deep scrambling experiments where circuit depth is the primary
fidelity bottleneck. The practical limit of 100 shots per circuit on AQT
(reflecting cost constraints) introduces statistical uncertainty
$\sigma \approx 0.10$ that limits quantitative precision. Future experiments with
${\sim}1000$ shots would reduce this to $\sigma \approx 0.03$, sufficient for
quantitative validation.

\subsection{System-Size Scaling and the $\Omega$ Flow}

The algebraic scaling $\Omega_{\rm late}(N) \approx 1/(1+AN^c)$ with $c \approx
1$ can be interpreted as a flow of the accessibility parameter $\Omega$ with the
system size $N$ as the renormalization scale. To our knowledge, this is the first
systematic analysis of the $N$-dependence of the OTOC saturation value as a flow;
Huang, Brand\~ao, and Zhang~\cite{HuangFiniteSize2019} established that this
quantity scales as $O(1/\text{poly}(N))$ but did not analyze the specific
functional form or its interpretation as a coupling flow.

The effective $\beta$-function $\beta_{\rm eff}(\Omega)$ describes how quickly
$\Omega$ flows from its UV value ($\Omega \to 1$ at small $N$) to its IR fixed
point ($\Omega \to 0$ at large $N$). The coefficient $c$ in
$\beta = -c\,\Omega(1-\Omega)$ (Hip.~A) or $\alpha$ in $\beta = -\alpha\,\Omega$
(Hip.~C) quantifies the speed of this flow. The observation
$c_{\rm SYK} \approx 3.9 \gg c_{\rm KI} \approx 1$ is consistent with SYK
saturating the Maldacena-Shenker-Stanford chaos bound~\cite{Maldacena2016},
indicating faster information delocalization per unit increase in system scale.
We note that our SYK-inspired model uses random ZZ all-to-all couplings rather
than the standard fermionic SYK interaction~\cite{SachdevYe1993,Kitaev2015};
the comparison to the MSS bound is therefore qualitative, reflecting the shared
feature of all-to-all connectivity rather than an exact correspondence.

A non-trivial property of the extracted $\beta_{\rm eff}(\Omega)$ is its
\textit{autonomy}: the function depends only on the current value of $\Omega$,
and not on the system size $N$ at which it is measured. This is a testable
consequence of the analysis, not an assumption. If $\beta_{\rm eff}$ depended on
$N$ or on the observer's position along the flow, data points extracted from
different system sizes would not collapse onto a single curve in the
$\beta$--$\Omega$ plane. The observed collapse across $N = 4$--$18$ (KI chaotic)
and $N = 4$--$12$ (SYK) constitutes direct empirical evidence that the flow is
autonomous: the rate of change of information accessibility is determined entirely
by the current state of accessibility, independently of the microscopic scale at
which the system is probed. This is the defining property of a Wilsonian
renormalization group flow, and to our knowledge constitutes the first empirical
characterization of an autonomous accessibility flow extracted from quantum
hardware.

% ============================================================
\section{Conclusion}
% ============================================================

We have characterized the spatial structure of quantum scrambling via
$\Omega_{\rm late}(r)$ profiles across five models spanning three dynamical
classes, validated the $\Omega$ diagnostic across three quantum hardware platforms
and two gate technologies, and analyzed the system-size scaling
$\Omega_{\rm late}(N)$ in the chaotic regime.

Spatial profiles reveal qualitatively distinct patterns: monotonically increasing
(KI chaotic, butterfly velocity signature), U-shaped (KI mixing, novel), and flat
(SYK, all-to-all), with integrable and dual-unitary systems uniformly at
$\Omega = 1$. Cross-platform comparison establishes ion trap hardware as superior
for deep scrambling circuits ($15\times$ lower error at $d=3$), providing the
first systematic benchmark across gate technologies for scrambling experiments.

System-size scaling analysis favors an algebraic form $\Omega_{\rm late} =
1/(1+AN^c)$ with $c \approx 1$ over a pure power-law ($\Delta{\rm AIC} = -3.30$),
with a markedly different coefficient $c \approx 3.9$ for the SYK model. The
effective $\beta$-function $d\Omega/d(\ln N)$ provides a new observable that
characterizes the speed of the accessibility flow with system scale, complementing
existing rate-based scrambling diagnostics. Furthermore, the effective
$\beta$-function is found to be autonomous---depending only on the current value
of $\Omega$ and not on system size---consistent with the defining property of a
Wilsonian renormalization group flow.

Full simulation code, raw data, circuit definitions, and analysis scripts are
available at~\cite{PerezEugenio2026a}. A supplemental document with complete
$\Omega(r)$ profiles for all models and system sizes, full cross-platform data
tables, and fit diagnostics accompanies this paper.

\begin{acknowledgments}
We acknowledge the use of IBM Quantum services, Amazon Web Services Braket
(Rigetti Ankaa-3 and AQT IBEX Q1 access). The author used AI-assisted tools for manuscript preparation and logical consistency checks. The views expressed are those of the
author and do not reflect the official policy or position of IBM, Amazon Web
Services, Rigetti Computing, or Alpine Quantum Technologies.
\end{acknowledgments}

\bibliography{kaelion_v8}

\end{document}
